% ============================================================
%  PLAN DE CONDUITE DU CHANGEMENT — Phase 2 : Planification
%  Time'Eats — Cellule PMO
%  Projet : Refonte Web
% ============================================================
\documentclass[11pt,a4paper,oneside]{report}
\usepackage{../../style/consulting}
\hypersetup{pdftitle={Plan de Conduite du Changement — Refonte Web — Time'Eats PMO}}
\pagestyle{fancy}\fancyhf{}
\renewcommand{\headrulewidth}{0pt}
\fancyhead[L]{\begin{tikzpicture}[remember picture,overlay]
  \draw[cnNavy,line width=0.5pt](0,0)--(\textwidth,0);
\end{tikzpicture}\vspace{2pt}\timeeatslogo\hspace{4pt}\small\textcolor{cnSlate}{Time'Eats \textbf{·} Conduite du Changement --- \textit{Refonte Web}}}
\fancyhead[R]{\small\textcolor{cnSlate}{\thepage}}
\fancyfoot[C]{\tiny\textcolor{cnSlate}{Document confidentiel --- Cellule PMO --- CESI MAALSI 2026}}
\begin{document}\small

\begin{center}
{\LARGE\bfseries\color{cnNavy} Plan de Conduite du Changement}\\[4pt]
\textcolor{cnOrange}{\rule{10cm}{1pt}}\\[4pt]
{\large\color{cnSlate} Phase 2 --- Planification \quad|\quad Projet : \textbf{Refonte Web}}
\end{center}

\vspace{6pt}

\begin{decision}
La refonte web modifie en profondeur le \textbf{workflow éditorial} (DSI $\rightarrow$ Marketing autonome via CMS) et introduit un \textbf{espace client authentifié}. Sans accompagnement, le risque d'usage faible du CMS et de retour aux anciens canaux (tickets DSI) est élevé. Ce plan structure l'adoption autour du modèle \textbf{ADKAR} (Awareness, Desire, Knowledge, Ability, Reinforcement).
\end{decision}

\section*{1 --- Analyse de l'impact sur les utilisateurs}

\renewcommand{\arraystretch}{1.4}
\begin{tabularx}{\textwidth}{L{4cm} L{2.5cm} Y C{2cm}}
\toprule
\rowcolor{cnNavy}\th{Population} & \th{Effectif} & \th{Impact du changement} & \th{Résistance attendue} \\
\midrule
Équipe Marketing / Communication (contributeurs CMS) & 12 personnes & \textbf{Très fort} : passage d'un workflow ticket DSI (5 j) à un CMS autonome (30 min). Nouveaux process de validation, calendrier éditorial. & \cellcolor{cnOrange!20}\textbf{Moyen} \\
\rowcolor{cnIce}
Équipe DSI Ops (run) & 5 personnes & \textbf{Fort} : nouvelle stack à exploiter (Next.js, Node, Azure), monitoring APM, fin du run Drupal historique. & \cellcolor{cnOrange!20}\textbf{Moyen} \\
Équipe Service Client (N1) & 8 personnes & \textbf{Moyen} : nouvelles FAQ, nouveau parcours espace client à diagnostiquer, formation à dispenser. & \cellcolor{cnGreen!20}\textbf{Faible} \\
\rowcolor{cnIce}
Équipes commerciales / partenariats & 15 personnes & \textbf{Moyen} : capacité d'auto-consommation de la nouvelle API par leurs outils, fin des exports manuels. & \cellcolor{cnGreen!20}\textbf{Faible} \\
Visiteurs publics / clients & N/A (externe) & \textbf{Léger} : nouvelle UX, redirections SEO, espace client (création de compte). & \cellcolor{cnGreen!20}\textbf{Faible} \\
\bottomrule
\end{tabularx}

\section*{2 --- Démarche ADKAR}

\renewcommand{\arraystretch}{1.4}
\begin{tabularx}{\textwidth}{L{2.5cm} L{3.5cm} Y C{2cm}}
\toprule
\rowcolor{cnNavy}\th{Étape ADKAR} & \th{Objectif} & \th{Actions concrètes} & \th{Timing} \\
\midrule
\textbf{A}wareness (Sensibilisation) & Comprendre \textit{pourquoi} la refonte & Newsletter de lancement, vidéo DG (2 min), affiches dans les bureaux, présentation CODIR & Avril--Mai 2026 (M1) \\
\rowcolor{cnIce}
\textbf{D}esire (Adhésion) & \textit{Vouloir} adopter le changement & Workshops co-construction maquettes UX avec contributeurs CMS, démo des bénéfices (gain de temps 5j $\rightarrow$ 30 min) & Mai--Juin 2026 (M2) \\
\textbf{K}nowledge (Savoir faire) & Savoir \textit{comment} utiliser le CMS & 4 sessions de formation CMS (3h chacune), tutoriels vidéo, guide PDF, sandbox d'entraînement & Août 2026 (M5) \\
\rowcolor{cnIce}
\textbf{A}bility (Capacité d'agir) & Réussir réellement à publier & Accompagnement individuel des Key Users en doublon les 2 premières semaines, hotline CMS dédiée & Septembre 2026 (M6) \\
\textbf{R}einforcement (Ancrage) & Pérenniser l'usage & Suivi mensuel KPI d'usage CMS, célébration des 100 premières publications, partage des bonnes pratiques en café-projet trimestriel & Oct.\ 2026 -- Janv.\ 2027 (M7--M10) \\
\bottomrule
\end{tabularx}

\section*{3 --- Réseau de Key Users}

\renewcommand{\arraystretch}{1.3}
\begin{tabularx}{\textwidth}{L{4cm} L{3cm} Y}
\toprule
\rowcolor{cnNavy}\th{Key User} & \th{Service} & \th{Rôle dans le déploiement} \\
\midrule
Mme H             & Marketing Contenu     & Référent CMS principal, animation des sessions Q\&R, relais bonnes pratiques \\
\rowcolor{cnIce}
Mme I              & Communication        & Référent ton de marque + SEO éditorial, formation des nouveaux contributeurs \\
M. J              & Commerce / partenariats & Référent contenus B2B et fiches partenaires \\
\rowcolor{cnIce}
Mme K               & Service Client       & Référent FAQ, parcours support, base de connaissances \\
M. D              & DSI Ops              & Référent run technique, monitoring, supervision \\
\bottomrule
\end{tabularx}

\section*{4 --- Plan d'actions opérationnel}

\renewcommand{\arraystretch}{1.3}
\begin{tabularx}{\textwidth}{C{0.6cm} L{4cm} Y C{2.5cm} C{2cm}}
\toprule
\rowcolor{cnNavy}\th{\#} & \th{Action} & \th{Description} & \th{Responsable} & \th{Timing} \\
\midrule
1 & Annonce projet              & Newsletter + vidéo DG + affichage & DG + CP + Marketing & 22/04/2026 \\
\rowcolor{cnIce}
2 & Workshop co-construction UX & 2 demi-journées avec contributeurs CMS & UX + Marketing      & 12--13/05/2026 \\
3 & Désignation Key Users       & Lettre de mission, planning dédié & CP + Marketing      & 27/05/2026 \\
\rowcolor{cnIce}
4 & Newsletter J1 (maquettes)   & Présentation des nouvelles maquettes & CP + Marketing      & 28/05/2026 \\
5 & Démo MVP (J3)              & Live démo en CODIR + tout Marketing & CP + ESN            & 31/07/2026 \\
\rowcolor{cnIce}
6 & Formations CMS (4 sessions) & 3h $\times$ 4 groupes = 12 contributeurs & CP + ESN + Mme H & 05--26/08/2026 \\
7 & Sandbox d'entraînement      & Environnement de test isolé pour les Key Users & DevOps + ESN        & 01--30/08/2026 \\
\rowcolor{cnIce}
8 & Documentation utilisateur   & Guide CMS PDF + 6 tutoriels vidéo + FAQ & CP + Marketing      & 25/08/2026 \\
9 & Hotline CMS Hypercare        & Permanence Teams + email dédié, J5+30 & DSI Ops + Key Users & 30/09--30/10/2026 \\
\rowcolor{cnIce}
10 & Enquête d'adoption           & Questionnaire NPS interne et logs d'usage & PMO + CP            & 30/10/2026 \\
11 & Café-projet trimestriel     & Bilan adoption + bonnes pratiques & PMO                 & 15/12/2026 \\
\bottomrule
\end{tabularx}

\section*{5 --- Indicateurs d'adoption}

\renewcommand{\arraystretch}{1.3}
\begin{tabularx}{\textwidth}{L{5cm} C{2cm} Y}
\toprule
\rowcolor{cnNavy}\th{Indicateur} & \th{Cible} & \th{Modalité de mesure} \\
\midrule
Taux de participation aux formations & $\geq$ 90\,\% & Feuilles d'émargement, plateforme LMS \\
\rowcolor{cnIce}
Satisfaction formation (post-session) & $\geq$ 80\,\% & Questionnaire en sortie de session \\
Publications réalisées dans le CMS à M+1 & $\geq$ 50 contenus & Logs CMS Strapi \\
\rowcolor{cnIce}
\% de publications CMS sans ticket DSI à M+3 & $\geq$ 90\,\% & Comparaison logs CMS vs.\ tickets ServiceNow \\
Temps moyen de publication d'un contenu & $<$ 30 min & Mesure CMS (création $\rightarrow$ publication) \\
\rowcolor{cnIce}
Incidents critiques CMS à J+30 & $\leq$ 5 & Log anomalies + tickets N1 \\
Score NPS interne contributeurs       & $\geq$ 30 & Enquête J5+30 \\
\rowcolor{cnIce}
Visites de la sandbox d'entraînement   & $\geq$ 2/contributeur & Logs CMS sandbox \\
\bottomrule
\end{tabularx}

\begin{reco}
La conduite du changement \textbf{ne s'arrête pas au Go-Live}. La phase Reinforcement (M7--M10) est cruciale : sans suivi des KPI d'usage à M+3, le risque de retour aux anciens canaux (tickets DSI) est réel. Le PMO inscrit le suivi mensuel d'adoption au cockpit DSI pendant les 6 mois post-Go-Live.
\end{reco}

\end{document}
